\section{Тестирование и отладка}
\subsection{Проверки Prolog}
В SWI-Prolog 10.0.2 выполнены 13 тестов plunit. Проверены положительный факт,
отсутствующая игра, игры Valve, сюжетные одиночные игры, сочетание И/ИЛИ/НЕ,
дружелюбный мультиплеер, множество напряжённых игр, сюжетные инди,
правило новичка, спокойные одиночные игры и десять пар игра--студия.
Ещё два теста демонстрируют false с отсечением и пять ответов без него.
Все 13 тестов прошли. Отдельный запуск десяти демонстрационных запросов
сохранён в evidence.json и использован для листингов отчёта.

\subsection{Проверки OWL и SWRL}
Использованы RDFLib 7.6.0, Owlready2 0.51 и встроенный в Owlready2 HermiT.
Это автоматический запуск того же семейства reasoner'ов, которое доступно
в Protégé; он не подменён простым обходом RDF-троек. Java взята из установленного Protégé.
После классификации наборы экземпляров шести определяемых классов и SWRL-класса
сравнены с ответами соответствующих предикатов Prolog.

\begin{table}[H]\centering\small
\caption{Проверки онтологии и фактический результат}
\begin{tabular}{P{8.7cm}P{5.5cm}}\toprule
Проверка & Результат\\\midrule
Изоморфность Turtle и RDF/XML & Графы совпадают\\
Консистентность и выполнимость именованных классов & Противоречий нет\\
Семь наборов выведенных типов против Prolog & Все совпали\\
Обратная связь developedGame для Valve & portal\_2, counter\_strike\_2\\
SoloStoryGame как подкласс SinglePlayerGame & Вывод подтверждён\\
SPARQL для сюжета без записанной связи с Valve & 6 игр\\
Minecraft одновременно Competitive и NonCompetitive & Обнаружена неконсистентность\\
Число 42 вместо строкового title & Обнаружена неконсистентность\\
Два различных title при exactly 1 & Обнаружена неконсистентность\\\bottomrule
\end{tabular}\end{table}

Негативные изменения вносятся в отдельные миры Owlready2 в памяти.
Файл исходной онтологии ими не изменяется. Это даёт возможность отличить
проверку корректного примера от проверки способности ограничений обнаруживать ошибки.

\subsection{Проверки диалога и DSS}
В Python unittest выполнены семь тестовых методов с параметризованными
случаями. Восемь интеграционных сценариев вызывают реальный SWI-Prolog.
Дополнительно проверены шесть неверных входных строк, четыре неверных набора
аргументов программного интерфейса, обе формы ввода (с префиксом и без него),
синонимы и дубликаты, отсутствие swipl,
повтор вопросов после ошибки и корректное завершение по EOF.
Все тестовые методы прошли; внешних сетевых сервисов приложение не требует.

\begin{lstlisting}
swipl -q -s lab_1/test_games.pl -g run_tests -t halt
python3 -m unittest discover -s lab_2 -v
python3 lab_1/part_2_ontology/verify_ontology.py
python3 reports/collect_evidence.py
\end{lstlisting}
Команды выполняются из module\_1. Для третьей команды нужны зависимости
из requirements-checks.txt и Java; для DSS достаточно стандартной библиотеки
Python и установленного SWI-Prolog. Пути к БЗ вычисляются относительно файлов
программы, поэтому запуск не зависит от текущей рабочей папки.

\subsection{Исправления по результатам аудита}
Исходные факты и семь правил первой работы соответствовали количественным требованиям.
Дополнены автоматические проверки, свойство данных title, кардинальности,
явный третий уровень иерархии и отдельно тестируемое SWRL-правило.
Уточнены документация по открытому миру и трактовка StudioGame.
Дубликат при дизъюнкции устранён на уровне выборки, а не изменением фактов.
Перед публикацией проверена обновлённая версия трёх файлов программы:
games.pl, bridge.pl и recommender.py. Описание необязательного префикса
ввода и коротких консольных сообщений приведено в соответствие с кодом;
тест ошибочного ввода RPG заменён на неизвестный RGB, поскольку RPG
без префикса теперь является допустимым вводом.
